\section{Conclusion}
\label{sec:conclusion}

We composed a learned biometric feature extractor with fuzzy-labelled PSI, replacing a
homomorphic matching layer with a public, training-free Super-Bit bridge. The engineering result
is favourable: 16-byte server-side templates against 1.67\,KB, no 117\,MB rotation key, no
slot-packing ceiling, and $1.88\%$ EER (95\% CI $[1.56, 2.20]$) against a $0.33\%$
floating-point ceiling.

The substantive result is a measurement. The error analysis that fixes FLPSI's parameters assumes
impostor codes are uniform. Codes from a 16-dimensional embedding are not: they are balanced ---
so the assumed \emph{mean} is exactly right and no first-moment check fires --- while carrying
$3.6\times$ the assumed standard deviation. At the published operating point this makes the
per-record false-accept rate $1{,}550\times$ larger than predicted, and no setting of the
sub-sampling parameters repairs it, because distinct fingers collide on identical codes at a rate
that lower-bounds every parameterisation.

The mechanism generalises beyond this system. An LSH bridge is valuable precisely because it is
faithful, and that faithfulness transmits the embedding's impostor tail into the code intact.
Any construction that matches learned codes under a threshold predicate inherits the same
problem, and can detect it cheaply: measure the impostor Hamming distribution rather than
assuming it, and count the zero-distance pairs. We recommend both as standard practice, and we
report multi-finger fusion as the only remedy we found that moves the bound --- four orders of
magnitude, where retuning moves none.
